\section{Introduction}
\label{sec:introduction}

\textbf{Infrastructure as Code (IaC)} describes servers, networks and cloud services in versioned source
files instead of manual configuration, and has become a research area of its own in Software
Engineering~\cite{rahman2019iacmapping}. In \textbf{Terraform}, the most widely used IaC tool, a
configuration is a set of \emph{blocks} written in HCL: resources, data sources, module calls,
input variables, outputs and locals. Blocks refer to each other by name, so a change rarely stays
local. Renaming an input variable breaks every block that still reads it and every module call that
still passes it (Listing~\ref{lst:example}); removing a resource breaks the outputs that expose it;
adding a bucket usually comes with its access policy and versioning blocks. Finding everything an
edit affects is the classic maintenance activity of \textbf{change impact
analysis}~\cite{bohner1996cia}. In IaC, a missed co-change surfaces late, when the configuration is
validated or planned, or never, when copies of a setting silently drift apart.

\begin{lstlisting}[style=hcl, caption={\small A variable rename inside a module and the two blocks it breaks.}, label={lst:example}, float=t]
# modules/storage/variables.tf
variable "bucket_name" { type = string }    # renamed from "name"

# modules/storage/main.tf
resource "aws_s3_bucket" "logs" {
  bucket = var.name                          # still reads the old name
}

# main.tf (root module)
module "storage" {
  source = "./modules/storage"
  name   = "app-logs"                        # still passes the old input
}
\end{lstlisting}

Several signals point to the blocks an edit affects, but none suffices alone. References are exact
but do not link copies of the same configuration kept in different directories. The commit
history reveals couplings that no reference shows~\cite{zimmermann2004rose}, but it is empty for new
blocks. Names and contents connect similar blocks, but also unrelated ones. How much each signal
should be trusted is not known beforehand and differs between them, which makes the combination a
natural target for machine learning: the task is to rank candidate blocks, the experience is the
commit history of many repositories, and performance is how high the blocks that really changed
are ranked~\cite{mitchell1997ml}. Learning the weights of the combination means searching a
continuous space for the best value of a ranking metric that has no gradient, a problem suited to
metaheuristic search and hence to search-based software engineering
(SBSE)~\cite{harman2001sbse}. The closest prior work, \textbf{GA-IaCRec}~\cite{bessghaier2026iacrec}, uses a
genetic algorithm to recommend co-changing \emph{files} for Ansible and Puppet, with the weights of
its two heuristics fixed by hand.

This report presents \tool, which learns this combination for Terraform at the level of
individual blocks, and answers 2 research questions:
\begin{itemize}
  \item \textbf{RQ1}: Is a genetic algorithm a feasible learner for block-level co-change
  prediction, i.e.\ does it learn from the labels, converge reliably, and find better weights than
  random search with the same budget?
  \item \textbf{RQ2}: Does a learned combination of signals rank co-changes better than mining the
  history alone (ROSE~\cite{zimmermann2004rose}) and than an equal-weight combination?
\end{itemize}

The contributions of this study are as follows:
\begin{enumerate}
    \item Block-level co-change recommendation for Terraform that also predicts
        the \emph{types} of blocks to add.
    \item A dataset of 53 open-source repositories mined commit by
        commit without information leakage from the future, plus 8 held-out repositories of a
        different kind.
    \item 12 signals in 3 families combined
        by weights learned with a self-implemented Genetic Algorithm;
    \item An evaluation on
        repositories never seen during training, together with the \tool{} command-line tool that applies
        the learned model to any local Terraform repository.
\end{enumerate}